\subsection{Las dieciséis contracciones del perfil}\label{sec:contracciones}
Denotamos por \(\langle g\rangle_i=\int_{[0,1]^2}\nu_1(i)g(i)\,di\) la media de área y
por \(\langle\cdot\rangle_c\) la misma media en el centro. En la fibra \(\lambda\)
de \(c\), \(d\omega_\lambda(i)=\eta_\lambda\nu_1(i)\,dy\), con \(y\) la coordenada libre
de fila o columna y, en las diagonales, la coordenada de fila. Para una función \(g\),
\(\int_F g=\sum_{\lambda\in \familias(c)}\int_{i\in\lambda}g(i,\lambda)\,d\omega_\lambda(i)\).
La medida \(d\rho\) de una línea libre es \(dy\) en
cada familia axial y \(U_f\,dy/(2x)\) en cada hoja de una familia diagonal.
El factor \(1/2\) ya realiza la media en la hoja de esa línea libre. Los demás índices
enteros libres aportan sus propias clases de paridad, según su parametrización. En una fibra
axial el desplazamiento tiene dos paridades de igual peso. En una diagonal el desplazamiento
mantiene el color. Un cruce diagonal--antidiagonal sólo existe en hojas compatibles.

\paragraph{Qué disponibilidad lleva una intersección.}
En una configuración fijamos las líneas de sus casillas marcadas y las líneas que se integran
con \(d\rho\). Sea \(\mathcal A\) ese conjunto. A cada otra línea \(\ell\) de una configuración
finita se le asigna una variable de Bernoulli independiente \(X_\ell\) de media \(U_\ell\).
En \(\mathcal A\) se pone \(X_\ell=1\). Sean \(\familias(z)\) las cuatro líneas de \(z\),
\(Y_z=\prod_{\ell\in \familias(z)}X_\ell\) y \(\pesodirecto(z)=W(z)Y_z\). Un cruce ausente da
 \(\pesodirecto=0\).
Así el peso de varias casillas directas \(Z\) se obtiene de
\[
 \psi(Z)=\mathbb E\prod_{z\in Z}Y_z
 =\prod_{\ell\in(\bigcup_{z\in Z}\familias(z))\setminus\mathcal A}U_\ell .
\]
El mismo indicador se usa en todo el monomio. Para el núcleo \(L_0\) y la corrección
\(\theta\) de la sección~\ref{sec:profile-data}, la respuesta de dos líneas es
\[
 E(a,b)=
 \begin{cases}
 (L_0)_{aa}+\theta_a,&a=b,\\
 (L_0)_{ab}-\pesodirecto(a\cap b),&a\ne b .
 \end{cases}
\]
Una casilla marcada aporta \(\pesodirecto(i)=W(i)\). Abreviamos
\(z_f=\sum_{j\in \familias(c)}E(f,j)\), \(T_i=\sum_{a\in \familias(i)}z_a\),
\(t_j=z_j\) para \(j\in \familias(c)\), y \(w(i)=\sum_{\{a,b\}\subset \familias(i)}(L_0)_{ab}\).
En particular \(t_j\) y \(w(i)\) en una casilla marcada son deterministas.
Para \(i\) que sólo comparte \(\lambda\) con \(c\), sea
\(\sigma_i=\sum_{a\in \familias(i)\setminus\{\lambda\}}z_a\).

\paragraph{Cubo, energía y corrección de fila.}
Las contribuciones del cubo de la acción son
\begin{align*}
 f2E&=\left\langle-\tfrac16\mathbb E T_i^3+\tfrac32 W(i)\mathbb E T_i^2\right\rangle_{i,c},\\
 f2z&=-\tfrac16\left\langle\int_{D\cup A}(u_f/U_f)\,\mathbb E z_f^3\,d\rho(f)\right\rangle_c,\\
 f2F&=\left\langle\int_F W(i)^2\right\rangle_c .
\end{align*}
En \(f2E\) se marcan \(i,c\), independientes en área. En \(f2z\) se marcan \(c,f\).
Para verificar los momentos en estos dos casos, sumemos sobre los pares de líneas que definen
\(T_i\) o \(z_f\). Sean \(A_1,A_2,A_3\) las sumas respectivas de \(L_0-K\),
\(K(W_{ab}-K)\) y \(-K(W_{ab}-K)(W_{ab}-2K)\), con \(W_{ab}=W(a\cap b)\) o cero.
En posición general las transversales de cruces distintos son distintas. El segundo momento es
\(A_1^2+A_2\) y el tercero \(A_1^3+3A_1A_2+A_3\), por los momentos centrados de cada Bernoulli.
El factor \(u_f/U_f\) es la vacancia del punto estacionario en la línea retenida.

Con \(I_2=\langle W(c)^2\rangle_c\) y \(\Xi=\langle W(c)w(c)\rangle_c\), se obtiene
\[
 G_{\rm cof}=-\tfrac12-3I_2,\qquad f4=-9I_2+\Xi,\qquad f5c=-3\Xi .
\]
En efecto \(mf_c\) tiende al factor \(w(c)-6W(c)\), y \(1-t_{cc}=-3+O_b(m^{-1})\).
Las otras contracciones del centro y de la fibra son
\begin{align*}
 f3F&=-\tfrac12\left\langle\int_F\mathbb E\sigma_i^2\right\rangle_c,\\
 f3c&=\tfrac12\left\langle\sum_{j\in \familias(c)}t_j^2
 -\sum_{j,h\in \familias(c)}[W(c)(1-\delta_{jh})+\Pi_{jh}]
                         [E(h,j)+\Pi_{hj}]\right\rangle_c,\\
 f5F&=-\left\langle\int_F w(i)\mathbb E T_i\right\rangle_c,\\
 f7&=\tfrac12\left\langle\left[-w(c)+\int_F w(i)
                  +\left\langle w(i)\mathbb E T_i\right\rangle_i\right]^2\right\rangle_c .
\end{align*}
Aquí \(\Pi_{jh}=v_jv_h/2\) es el núcleo del perfil. El primer sumando de \(f3c\)
procede de \(t^{\mathsf T}D_0t\) y el segundo de \(E_R\).
En \(f3F\), \(i,c\) están marcadas en la misma fibra. Sean \(Z\) los cruces existentes
\(a\cap j\), con \(a\in \familias(i)\setminus\{\lambda\}\) y \(j\in \familias(c)\setminus\{\lambda\}\),
de familias distintas. En posición general son casillas distintas. Escribiendo
\(\psi_z=\psi(\{z\})\), el signo negativo de cada parte directa desaparece en la covarianza,
\begin{align*}
 \mathbb E\sigma_i^2={}&(\mathbb E\sigma_i)^2+\sum_{z\in Z}W(z)^2\psi_z(1-\psi_z)\\
 &+\sum_{\substack{z,z'\in Z\\z\ne z'}}W(z)W(z')
       [\psi(\{z,z'\})-\psi_z\psi_{z'}].
\end{align*}
Por ejemplo, en coordenadas relativas a \(c=(0,0)\), si \(i=(0,s)\), los cruces
\((s,s)\) y \((s,0)\) comparten la fila \(s\). Cuando existen dentro del tablero, su peso
conjunto es \(\psi_z\psi_{z'}/U_R\). En \(f7\) las dos copias de la suma tienen índices
libres distintos. Tras fijar \(c\), sus coincidencias accidentales tienen medida nula.
Así se cuadran las integrales de la respuesta centrada antes de promediar en \(c\).

\paragraph{Triángulos de líneas en la traza cúbica.}
Tres líneas del producto \(S_*^3\) concurren en una casilla borrada o cortan tres casillas
distintas que completan un cuadrado con \(c\). En el primer caso sus seis órdenes cancelan
el divisor \(6\), y el resultado es otra copia de \(f2F\).
Para el segundo, tomemos la base \(((1,0),(0,1))\) o \(((1,1),(1,-1))\) y multipliquemos
cada vector por un signo independiente, obteniendo \(v,w\). Sean
\(Z=\{c,c+hv,c+hw,c+h(v+w)\}\), \(h>0\), y \(n_{RC}\) el número de líneas axiales de \(Z\).
Las tres coordenadas libres son las de \(c\) y el paso primitivo \(h\). El resultado es
\[
 C_{3,\rm cuad}=\sum_{v,w}\int_{\{(c,h):\,Z\subset[0,1]^2\}}
  x^{n_{RC}-3}\prod_{z\in Z}W(z)
  \prod_{\ell\in(\bigcup_{z\in Z}\familias(z))\cap(D\cup A)}U_\ell\,dc\,dh .
\]
Los cuatro factores \(W\) incluyen el peso del centro. En la cuadratura el lado cero
lleva la mitad del valor límite del integrando. En la integral tiene medida cero.

\paragraph{Respuesta de segundo orden en el exterior.}
La fila \(f6\) requiere \(m^2q_i\), pues la respuesta de primer orden en \(y_i/p_i\)
se cancela con \(t_{ic}\). Para verificar la fuente de segundo orden, pongamos
\(X=(a_{i,I}^{\mathsf T})_{i\in E'}\), \(D_p=\operatorname{diag}(p_i)_{i\in E'}\).
Las ecuaciones estacionarias restadas dan \(y_i/p_i=(X\eta)_i+(B_1)_i\), donde
\(\eta_f=\gamma_f/h_f\) en las diagonales y las otras componentes son multiplicadores
de fila y columna. Con \(\theta_i=y_i/p_i\), la diferencia exacta de los logaritmos
de las ecuaciones de estacionariedad da
\[
 (B_1)_i=\theta_i-\log(1+\theta_i)-3\log(1-y_i/(1-p_i))
 +\sum_{f\in \familias(i)\cap(D\cup A)}[\log(1+\eta_f)-\eta_f].
\]
Las ecuaciones de masa dan
\(\mathcal M\eta=d_0-X^{\mathsf T}D_pB_1\). Por tanto
\[
 q=B_1+X\mathcal M^+(r-X^{\mathsf T}D_p B_1).
\]
Taylor en las casillas y las vacancias da \(m^2B_1=L+O_b(m^{-1})\), con
\[
 L_i=\tfrac12(mt_{ic})^2-\frac{3mp_i(mt_{ic})}{1-p_i}
       -\tfrac12\sum_{f\in \familias(i)\cap(D\cup A)}(mz_f^{\rm fis})^2,\qquad
 z^{\rm fis}=(Ra_c)_I .
\]
Aquí se han usado \(y_i/p_i=-t_{ic}+O_b(m^{-2})\) y
\(\gamma_f/h_f=-z_f^{\rm fis}+O_b(m^{-2})\).
En el complemento del modo nulo, la identidad \(S_0R=I-\Pi_m\) por bloques da
\(\|\mathcal M^+-R_{II}\|_\infty=O_b(m^{-1})\) en norma de suma por fila,
pues las líneas retenidas pierden masa \(O_b(m^{-1})\) y la inversa está acotada.
Así \(m^2q=L+XR_{II}(m^2r-X^{\mathsf T}D_pL)+O_b(m^{-1})\).

Expresamos la contracción en los mismos núcleos del perfil. Fijemos el centro \(c\) y sea
\(\Lambda(i)=T_i^2/2-3W(i)T_i-\tfrac12\sum_{f\in \familias(i)\cap(D\cup A)}z_f^2\).
Para una línea retenida \(u\), escribamos
\begin{align*}
 \omega(u)={}&-\sum_{j\in \familias(c)} \pesodirecto(u\cap j)
     \left[t_j+\sum_{a\in \familias(u\cap j)\setminus\{j\}}z_a\right]
      -\int_{i'\in u}\mathbb E\Lambda(i')\,d\omega_u(i'),\\
 q(i)={}&\Lambda(i)+\sum_{u\in \familias(i)}
 \left[\omega(u)+\int E(u,u')\omega(u')\,d\rho(u')\right],\\
 f6={}&-\left\langle w(i)\mathbb E q(i)\right\rangle_{i,c}.
\end{align*}
Se omiten los cruces inexistentes en la primera suma. El término de fibra integrado en
\(\omega\) es la media con \(i'\) marcada y \(c\) fijo. En los productos de la suma sobre
cruces se usa \(\psi\) para todo el monomio. Las dos partes de \(\omega\) provienen de
\(m^2r=-m^2(C_0t+U_0z)\) y de \(-X^{\mathsf T}D_pL\). La identidad y la parte integral
de \(R_{II}\) producen los dos términos entre corchetes de \(q(i)\). La media de área
lleva el signo de \(-m^2\sum_{E'}p_iq_iw_i\). El Lema~\ref{lem:f6} demuestra su
reordenación por trazas, incluido el recorrido de la figura~\ref{fig:rectangulos}.

\begin{lema}[Reordenación de la contribución \(f6\)]\label{lem:f6}
En la malla de lado \(s\), sean \(\mathcal I{}\) la incidencia casilla--línea de la
sección~\ref{sec:rectangulos}, \(\beta_e=\nu_1(e)/s^2\) y
\(N_v=\mathcal I{}^{\mathsf T}\operatorname{diag}(v)\mathcal I{}\). Los productos de vectores de casillas
son puntuales. Con \(K,L_0,\rho,w\) de la sección~\ref{sec:profile-data}, definamos
\begin{align*}
 D&=L_0-K,& V_{ij}&=K_{ij}(W_{ij}-K_{ij}),&
 \Psi&=\rho^{-1}\mathcal I{}^{\mathsf T}(\beta w),\\
 Z&=\Psi+D\rho\Psi,& h&=w-\mathcal I{}Z,& \alpha&=\beta h,\\
 a&=\mathcal I{}^{\mathsf T}\alpha,& b_\beta&=\mathcal I{}^{\mathsf T}\beta,&
 M&=\mathcal I{}D\mathcal I{}^{\mathsf T}.
\end{align*}
Aquí \(W_{ij}\) es el peso del cruce o cero si éste no existe.
Escribamos \(U(e,j)=\sum_{u\in \familias(e)} K_{uj}-3W(e)\) y
\(t(c,e;j)=M_{ce}+\theta_j+U(c,j)+U(e,j)\).
La cuadratura del perfil físico, con el valor genérico en coincidencias accidentales,
da exactamente \(f6_s=-E_1+E_2\), donde
\begin{align*}
 E_1={}&\tfrac12\operatorname{tr}(N_\beta D N_\alpha D)
       +\tfrac12 b_\beta^{\mathsf T}Va-3b_\beta^{\mathsf T}D\mathcal I{}^{\mathsf T}(\alpha W)\\
     &-\tfrac12\sum_{f\in D_{\rm obl}}a_f
         [(D N_\beta D)_{ff}+(Vb_\beta)_f],\\
 E_2={}&\sum_c\beta_c\sum_{j\in \familias(c)}\sum_{e\in j}
     \frac{\eta_j\nu_1(e)}s\,t(c,e;j)
     \sum_{u\in \familias(e)\setminus\{j\}}Z_u .
\end{align*}
En esta fórmula \(D_{\rm obl}\) reúne las dos familias oblicuas.
\end{lema}
\begin{proof}
Fijado \(c\), la respuesta de segundo orden consta de la fuente local \(\Lambda\) y
la fuente de líneas \(\omega\). Al sumar la respuesta contra \(\beta_e w(e)\), cada
\(u\in \familias(e)\) recibe \([\mathcal I{}^{\mathsf T}(\beta w)]_u=(\rho\Psi)_u\).
La parte integral de la respuesta añade \(\rho D\rho\Psi\) por simetría. La masa
total es, pues, \(\rho Z\). La fuente negativa de \(\omega\) integra \(\Lambda\) en
la fibra de \(u\). En la malla \(\rho_u\eta_u=1/s\), por lo que esa suma sustrae
\(\beta_e(\mathcal I{}Z)_e\) del peso de \(\Lambda(e;c)\). Queda la masa \(\alpha_e\).

Para dos marcas genéricas \(c,e\), la media de \(T_e\) es \(M_{ce}\).
Cada cruce directo tiene varianza \(V_{uj}\), pues su indicador al cuadrado es el propio
indicador. Los cruces distintos tienen transversales distintas en posición genérica.
Así \(\mathbb E T_e^2=M_{ce}^2+\sum_{u\in \familias(e),j\in \familias(c)}V_{uj}\).
Para \(f\in \familias(e)\) la misma cuenta da
\(\mathbb E z_f^2=(\sum_{j\in \familias(c)}D_{fj})^2+\sum_{j\in \familias(c)}V_{fj}\).
Al sumar \(\beta_c\alpha_e\mathbb E\Lambda(e;c)\), el cuadrado \(M_{ce}^2\)
da la traza de \(N_\beta D N_\alpha D\) por multiplicación de índices.
La varianza da \(b_\beta^{\mathsf T}Va\), y el momento lineal con \(W(e)\)
da \(b_\beta^{\mathsf T}D\mathcal I{}^{\mathsf T}(\alpha W)\).
Los cuadrados de \(z_f\) dan \((DN_\beta D)_{ff}+(Vb_\beta)_f\), uno por línea
oblicua de \(e\). Los coeficientes de \(\Lambda\) producen \(E_1\).
En \(N_v\) se incluyen las entradas diagonales \(\sum_{e\in u}v_e\), también cuando
el recorrido vuelve por la misma línea.

En la suma de cruces de \(\omega\), fijemos \(e=u\cap j\), con \(j\in \familias(c)\).
El peso directo por la medida de \(u\) es \(\rho_uK_{uj}=\eta_j\nu_1(e)/s\).
Condicionar el cruce a estar disponible marca entonces \(c,e\) en la fibra \(j\).
En la respuesta media \(M_{ce}\) se repone la diagonal \(\theta_j\).
Los tres cruces con \(j\) en cada extremo llevan ahora el peso cierto \(W\) de la
casilla marcada. Cambiar cada antiguo término \(-K\) a \(-W\) añade \(U(c,j)+U(e,j)\).
Queda \(t(c,e;j)\), multiplicado por \(Z_u\) para cada transversal \(u\ne j\) de \(e\).
La fuente de cruces tiene signo negativo en \(\omega\), y la definición de \(f6\)
tiene otro signo negativo. De ahí \(+E_2\) y \(-E_1\). Esta igualdad de sumas
finitas deja el error espacial de la extrapolación en el estado \status{N}.
\end{proof}

\paragraph{Las dos trazas del término mixto.}
Pongamos \(D=H_*D_*-S_*\). La ciclicidad de la traza da la identidad en el padre
\[
 Q_c=\tfrac12\sum_{u,v}(S_*)_{uv}(m^2D)_{vu}
       +\tfrac1{4m^2}\sum_{u,v}(m^2D)_{uv}(m^2D)_{vu}.
\]
Cada casilla de \(F\setminus\{c\}\) tiene tres líneas retenidas, que dan seis pares
ordenados en la primera suma. En la segunda se integran dos líneas libres.
Para expresar \(m^2D\) en el perfil escribamos \(\mathcal K(v,w)=E(v,w)+\Pi_{vw}\),
\(\mathcal V(v,j)=\pesodirecto(v\cap j)+\Pi_{vj}\), \(j\in \familias(c)\). Para líneas retenidas \(v,u\),
sea \(\mathcal Y(v,u)=-\pesodirecto(z)T_z\) si se cortan en \(z\in E'\), y cero en otro caso.
En esta fórmula \(T_z\) usa los mismos indicadores que \(\pesodirecto(z)\).
Finalmente, para \(w\ne u\) retenidas, sea \(\mathcal S(w,u)=-\pesodirecto(i')\) cuando
\(w\cap u=i'\in F\setminus\{c\}\), y cero en los demás pares. Definamos
\begin{align*}
 \mathcal B(v,u)&=\mathcal Y(v,u)+\Pi_{vu}
                    +\sum_{j\in \familias(c)}\mathcal V(v,j)\mathcal V(u,j),\\
 \mathcal E_2(v,u)&=\mathcal B(v,u)+\sum_w\mathcal K(v,w)\mathcal S(w,u)
                     +\int\mathcal K(v,w)\mathcal B(w,u)\,d\rho(w).
\end{align*}
Estas tres piezas proceden, con \(B_*=D_*-S_*\) y \(K_*=H_*-I\), de
\(m^2D=m^2B_*+(mK_*)(mS_*)+m^{-1}(mK_*)(m^2B_*)\).
En la suma discreta se toma cada \(i'=u\cap j\), \(j\in \familias(c)\), y cada una de
sus dos líneas transversales \(w\). Su disponibilidad está en \(\pesodirecto(i')\).
En la integral sí se marca \(w\), pues su disponibilidad ya está en \(d\rho(w)\).

Escribamos \(R(i)=\{(u,v):u,v\in \familias(i)\setminus\{\lambda\},\ u\ne v\}\) para una
casilla de la fibra \(\lambda\). Aplicando \(\psi\) a los monomios se obtienen
\begin{align*}
 L_1&=-\tfrac12\left\langle\int_F
          \sum_{(u,v)\in R(i)}\mathbb E\mathcal E_2(v,u)\right\rangle_c,\\
 L_2^{\rm int}&=\tfrac14\left\langle\iint
 \mathbb E\big[\mathcal E_2(u,v)\mathcal E_2(v,u)-\mathcal Y(u,v)\mathcal Y(v,u)\big]
                 \,d\rho(u)\,d\rho(v)\right\rangle_c .
\end{align*}
Las dos copias de una línea integrada internamente tienen variables libres separadas.
El bloque \(\mathcal Y\mathcal Y\) se sustrae antes de integrar. En el tablero cada
\(i\in E'\) aporta doce pares ordenados al bloque \(YY\), de modo que su contribución
es \(+3m^2\sum_{E'}y_i^2\), exactamente la opuesta a la de \(P_c\).
También se puede comprobar en el continuo. Para cada par de familias, la medida de los
cruces válidos por su disponibilidad es \(U_DU_A\,di\), con las hojas y el jacobiano del
cruce incluidos. Multiplicando por \(W(i)^2\) resulta \(\nu_1(i)W(i)\).

\paragraph{Cómo extraer \(\Delta_1\) y \(\Delta_2\).}
Las filas tabuladas son \(Q_1=L_1-\Delta_1\) y \(Q_2=L_2^{\rm int}-\Delta_2\).
Damos las correcciones por separado para poder reconstruir cada entrada. En todas las
sumas siguientes un cruce debe existir en el cuadrado y en la clase de paridad indicada.
Una igualdad de líneas se exige como identidad en los parámetros libres de esa clase.
Las coincidencias accidentales forman conjuntos de medida nula.

Para \(\Delta_1\), fijemos \(c,i,\lambda,(u,v)\in R(i)\) como en \(L_1\), y una
línea libre \(w\) integrada con \(d\rho(w)\). Pongamos \(z=v\cap w,\ p=w\cap u\).
Para cada \(a\in \familias(p)\setminus\{w,u\}\) y \(j\in \familias(c)\), pongamos \(t=a\cap j\).
Se marcan \(\familias(c)\cup \familias(i)\cup\{w\}\).
El monomio directo de \(\mathcal K(v,w)\mathcal Y(w,u)\) aporta la diferencia
\[
 \Delta_1=\tfrac12\left\langle\int_F\sum_{(u,v)\in R(i)}
 \int\sum_{a,j} W(z)W(p)W(t)
 \big[\psi(\{z,p,t\})-\psi(\{z\})\psi(\{p,t\})\big]\,d\rho(w)\right\rangle_c .
\]
Se toman los tres cruces distintos, con \(p\in E'\). La diferencia cuenta la transversal
común de \(z,t\) sólo una vez. Si es \(\ell\), el corchete vale
\(\psi(\{z\})\psi(\{p,t\})(U_\ell^{-1}-1)\).
Las tres partes directas llevan signo negativo en el producto de la primera traza
antes de aplicar su prefactor \(-1/2\).
Por ejemplo, en coordenadas relativas, para \(c=(0,0)\), \(i=(0,s)\), \(u=C_s\),
\(v=D_{-s}\), \(w=R_h\), \(a=A_{h+s}\), \(j=R_0\), resultan \(z=(h,h+s)\),
\(p=(h,s)\), \(t=(0,h+s)\). La columna \(h+s\) aparece en \(z\) y \(t\).

Para \(\Delta_2\), marquemos \(\familias(c)\cup\{u,v\}\), con \(u,v\) libres. Cuando existe el cruce,
el peso de sus transversales es
\[
 K_{ab}=W(a\cap b)\prod_{\ell\in \familias(a\cap b)\setminus\{a,b\}}U_\ell .
\]
Para cada \(j\in \familias(c)\) tomamos \(p_2=v\cap j\),
\(w_2\in \familias(p_2)\setminus\{v,j\}\), \(t_2=u\cap w_2\), \(k_2=K_{vj}K_{uw_2}\).
Los dos patrones posibles para los otros dos vértices dan
\[
\begin{array}{c|c|c|c}
 &p_1,\ t_1 &k_1&\text{multiplicidad }s\\ \hline
 \mathcal K\mathcal S\ \text{por}\ \mathcal K\mathcal S&
 u\cap j,\ v\cap w_1 & K_{uj}K_{vw_1}&1\\
 \mathcal Y\ \text{por}\ \mathcal K\mathcal S&
 v\cap u,\ a\cap j & K_{vu}K_{aj}&2
\end{array}
\]
En la primera fila \(w_1\) recorre \(\familias(p_1)\setminus\{u,j\}\). En la segunda \(a\)
recorre \(\familias(p_1)\setminus\{v,u\}\). En cada patrón se conservan cuatro vértices distintos.
Para cada pareja de cruces que define \(k_d\), \(d=1,2\), reunamos las transversales de
sus dos cruces, es decir, las líneas cuyos factores \(U\) aparecen en \(k_d\).
Si los dos conjuntos comparten una línea \(\ell\) fuera de las marcadas, se incluye el
sumando siguiente. Los cuatro cruces forman entonces un cuadrado y la transversal es única.
\[
 \Delta_2=\tfrac14\left\langle\iint\sum_{\text{patrones admitidos}}
          s\,k_1k_2(U_\ell^{-1}-1)\,d\rho(u)\,d\rho(v)\right\rangle_c .
\]
Se puede obtener la lista desarrollando los productos de \(\mathcal E_2\).
Un peso \(\pesodirecto\) aporta las cuatro coordenadas de su casilla menos las líneas ya marcadas.
Para comparar dos etiquetas de una misma familia se restan sus expresiones afines en
las coordenadas libres. Sólo una diferencia idénticamente nula aporta el factor
\(U_\ell^{-1}-1\). Para hacer la comprobación basta comparar el término constante y
los coeficientes de cada parámetro en las cuatro etiquetas \(r,q,r-q,r+q\). En
\(\Delta_1\) se obtienen 24 elecciones de familias y 80 al distinguir las clases admisibles.
En \(\Delta_2\) quedan 24 elecciones en cada fila de la tabla. Por ejemplo, en la primera,
con \(j=R_0,u=A_s,v=D_t,w_1=C_s,w_2=C_{-t}\), los dos vértices exteriores
\((s+t,s),(s+t,-t)\) comparten la fila \(s+t\). En la parte mixta, intercambiar las copias da el mismo
integrando con \(u,v\) permutadas, de donde procede el \(2\) de la tabla.
Finalmente,
\begin{align*}
 d(x)={}&G_{\rm cof}+f2E+f2z+2f2F+f3F+f3c+f4+f5c+f5F\\
        &+f6+f7+C_{3,\rm cuad}+Q_1+\Delta_1+Q_2+\Delta_2 .
\end{align*}

